\documentclass[10pt,a4paper]{amsart}
\usepackage[margin=19mm]{geometry}
\usepackage{amsmath,amssymb,mathtools}
\usepackage[scr=rsfs,cal=euler]{mathalfa}
\usepackage{xfrac}
\usepackage[unicode,hidelinks,bookmarks=false]{hyperref}
\newcommand{\ie}{i.e.\ }
\newcommand{\cf}{cf.\ }
\newcommand{\resp}{respectively\ }
\newcommand{\Subsection}{\subsection}
\providecommand{\nobreakdash}{\nobreak\hbox{-}}
\setlength{\emergencystretch}{2em}
\multlinegap0pt
\usepackage{amssymb}
\usepackage{xcolor}
\newtheorem{lemma}{Lemma}
\newcommand{\N}{\mathbb{N}}
\newcommand{\Q}{\mathbb{Q}}
\title{The spectrum of $(\xi\alpha^n)$ can be uncountable}
\author{H\.IKMET BURAK \"OZCAN}
\date{Source: arXiv:2609.07714v1 (CC BY 4.0)}
\begin{document}
\maketitle

\noindent\emph{Source and licence.} The spectrum of $(\xi\alpha^n)$ can be uncountable. Version arXiv:2609.07714v1 (CC BY 4.0), distributed by arXiv under the Creative Commons Attribution 4.0 International licence, http://creativecommons.org/licenses/by/4.0/, as recorded in the arXiv licence metadata for this exact version and shown on the abstract page https://arxiv.org/abs/2609.07714v1. Full licence notice: \texttt{licenses/arxiv-2609.07714-CC-BY-4.0.txt}.\par\smallskip

\noindent\textit{Editorial scope.} Counted unit: the article's lemma lem:cantor (source0.tex lines 166--168). The article's complete original proof of it is delivered with it (source0.tex lines 170--202, one \texttt{proof} environment of 33 lines). No mathematics is written, repaired or re-derived: the statement and the proof are the article's own lines, in the article's own notation, and no retained line is substituted or re-flowed.  Retained uncounted as the source-local context the proof needs: lines 154--156.  Nothing was omitted from any retained span, so the package carries no omission record.  No retained line contains a citation, so the wrapper carries no bibliography and no bibliography entry.  No retained line contains a figure environment, an image-inclusion command or drawing code, so no image is delivered and the package declares no asset dependency.  Editorial formatting only, in addition: an AMS \texttt{amsart} preamble that restates the article's own declarations and macros used by the retained lines, namely the environments \texttt{lemma} on separate counters, together with the standard packages those lines need.  The retained environments print in this document as Lemma 1 in order of appearance, whereas the article prints the same items with the numbering of its own full document.  The complete native source of the article as distributed in the arXiv e-print for this exact version is bundled unchanged as \texttt{sources/209652/source0.tex}.
\begin{lemma}\label{lem:tree}
For every $j\ge1$, the intervals $J_{2j}$ and $J_{2j+1}$ are disjoint subintervals of $J_j$. Consequently, we have $J_j\subseteq J_1\subset(0,1)$ for every $j\ge1$.
\end{lemma}

\begin{lemma}\label{lem:cantor}
There is an uncountable set $C\subset(0,1)$ such that all its elements are irrational and every $\theta\in C$ lies in $J_j$ for infinitely many $j$.
\end{lemma}

\begin{proof}
Let $\omega=(\omega_k)_{k\ge0}\in\{0,1\}^{\N}$. We define the nodes along a branch by setting $n_0(\omega)=1$ and by moving to the left or to the right child at each step, that is,
\[
n_{k+1}(\omega)=2n_k(\omega)+\omega_k .
\]
We also put $\sigma_k=+1$ if $\omega_k=1$, and $\sigma_k=-1$ if $\omega_k=0$. By the recursive definition of the centres, the centre of the $K$-th interval along this branch is the partial sum
\begin{equation}\label{eq:partial}
\theta_{n_K(\omega)}=\tfrac12+\sum_{k<K}\sigma_k\,\frac{r_{n_k(\omega)}}2\qquad(K\ge0).
\end{equation}
Since $n_k\ge1$ and $n_0=1$, we get $n_{k+1}\ge2n_k\ge n_k+1$, and hence by induction
\[
n_k\ge k+1\qquad\text{and}\qquad n_{K+i}\ge n_K+i .
\]
Recalling that $r_j=\kappa/R^j$, we deduce
\begin{equation}\label{eq:radii}
r_{n_k}\le\frac{\kappa}{R^{\,k+1}}\qquad\text{and}\qquad r_{n_{K+i}}\le\frac{\kappa}{R^{\,n_K+i}}=\frac{r_{n_K}}{R^{\,i}}.
\end{equation}
By the first bound in \eqref{eq:radii}, the series
\[
\theta(\omega):=\tfrac12+\sum_{k\ge0}\sigma_k\,\frac{r_{n_k(\omega)}}2
\]
converges absolutely. To see that $\theta(\omega)$ lies in every interval along its branch, we subtract \eqref{eq:partial} from this series, which leaves the tail $\theta(\omega)-\theta_{n_K(\omega)}=\sum_{i\ge0}\sigma_{K+i}\,r_{n_{K+i}}/2$. As $|\sigma_k|=1$, the second bound in \eqref{eq:radii} and $R/(R-1)\le\tfrac32$, which holds since $R\ge3$, give
\[
\bigl|\theta(\omega)-\theta_{n_K(\omega)}\bigr|\le\sum_{i\ge0}\frac{r_{n_{K+i}}}2\le\frac{r_{n_K}}2\sum_{i\ge0}\frac1{R^{\,i}}=\frac{r_{n_K}}2\cdot\frac{R}{R-1}\le\frac34\,r_{n_K}<r_{n_K}.
\]
Therefore, $\theta(\omega)\in J_{n_K(\omega)}$ for every $K\ge0$, and since the indices $n_K(\omega)$ are strictly increasing, $\theta(\omega)$ lies in infinitely many $J_j$.

Next, we show that the map $\omega\mapsto\theta(\omega)$ is injective. Let $\omega\ne\omega'$, and let $\ell$ be the first index where they differ. The two branches share the same nodes up to $\ell$, so that $n_k(\omega)=n_k(\omega')$ for $k\le\ell$, whereas $n_{\ell+1}(\omega)$ and $n_{\ell+1}(\omega')$ are the two distinct children $2n_\ell$ and $2n_\ell+1$. By Lemma~\ref{lem:tree}, the corresponding intervals are disjoint, and they contain $\theta(\omega)$ and $\theta(\omega')$ respectively. Hence, $\theta(\omega)\ne\theta(\omega')$, and the range of $\theta$ has the cardinality $2^{\aleph_0}$ of $\{0,1\}^{\N}$. Now, we set
\[
C:=\bigl\{\theta(\omega):\omega\in\{0,1\}^{\N}\bigr\}\setminus\Q.
\]
Since $\Q$ is countable, $C$ is uncountable. All of its elements are irrational, they lie in $J_1\subset(0,1)$ by Lemma~\ref{lem:tree}, and each of them lies in infinitely many $J_j$ by the first part of the proof.
\end{proof}

\end{document}
